% GPT-generated supplementary answers. Review before using in instruction.
% ANSWERS: f1
\begin{itemize}
\item \textbf{教材习题 1.4：}参见教材。
\item \textbf{数列单调性：}令 $u_n=(1+1/n)^n$。由均值不等式，$u_n$ 严格递增；由二项展开或 Bernoulli 不等式，$v_n=(1+1/n)^{n+1}$ 严格递减。两列同趋于 $e$。
\item \textbf{Cauchy--Schwarz 等号：}正确条件为向量 $(a_k)$ 与 $(b_k)$ 线性相关，包括其中一个零向量的情况。原稿用非零常数表示的写法漏掉了部分零向量情形。
\item \textbf{双曲函数反函数：}$\sinh:\mathbb R\to\mathbb R$ 严格递增，反函数为 $\operatorname{arsinh}y=\ln(y+\sqrt{1+y^2})$。$\cosh$ 在 $\mathbb R$ 上不是一一对应；限制到 $[0,\infty)$ 后，$\operatorname{arcosh}y=\ln(y+\sqrt{y^2-1})$，$y\ge1$。
\item \textbf{两道不动点题：}若 $f^2$ 的唯一不动点为 $a$，则 $f(a)$ 也被 $f^2$ 固定，故 $f(a)=a$；任何 $f$ 的不动点也是 $f^2$ 的不动点。若 $f^2$ 恰有 $a,b$ 两个不动点，则 $f$ 将集合 $\{a,b\}$ 映到自身，且在该集合上 $f^2=\mathrm{id}$，因而只能恒等或交换两点。
\item \textbf{集合原像：}$C\subseteq f^{-1}(f(C))$，$f(f^{-1}(D))=D\cap f(A)\subseteq D$。分别对所有子集取等号当且仅当 $f$ 单射、满射。
\end{itemize}

% ANSWERS: f2
\begin{itemize}
\item \textbf{教材章节题号：}参见教材。
\item \textbf{进制截断：}取 $k_n=\lfloor m^n x\rfloor$，有 $k_n\le m^n x<k_n+1$。由 $\lfloor mt\rfloor\ge m\lfloor t\rfloor$ 与 $\lfloor mt\rfloor\le m\lfloor t\rfloor+m-1$ 得 $x_n\uparrow$、$y_n\downarrow$；间距 $y_n-x_n=m^{-n}$。十进制时二者分别是下、上截断。
\item \textbf{有理数中的上确界：}若 $r\in\mathbb Q_{>0}$ 且 $r^2<2$，利用稠密性可在 $r$ 与 $\sqrt2$ 间取有理数，故 $r$ 非上界；若 $r^2>2$，可在 $\sqrt2$ 与 $r$ 间取更小的有理上界。$r^2=2$ 不可能，故该集合在 $\mathbb Q$ 中无上确界。
\item \textbf{底数 $q$ 的表示：}先用阿基米德性质确定唯一 $k$ 使 $q^k\le x<q^{k+1}$，再逐位取整数部分 $d_j$。每取一位后余数落在长度为 $q^{k-j}$ 的区间内，归纳得到题中不等式。
\item \textbf{稠密性、根与直径：}有理数稠密性加任一固定无理数的平移给出无理数稠密性；$x\mapsto x^n$ 在正半轴严格递增，利用上确界构造唯一正根；集合直径用 $|x-y|\le\sup S-\inf S$ 和逼近上下确界的两列证明。
\item \textbf{选做题：}有理指数运算法则可归结到统一分母后的整数指数法则。Dirichlet 逼近采用抽屉原理；若只有有限多个 $p/q$ 满足 $|\alpha-p/q|<q^{-2}$，对充分大的逼近参数应用 Dirichlet 定理即矛盾。
\end{itemize}

% ANSWERS: f3
\begin{itemize}
\item \textbf{教材章节题号：}参见教材。
\item \textbf{$3、4、7、8$ 班极限：}由于 $1\le k^{1/n}\le n^{1/n}$，第一极限为 $1$。第二题先减去 $\sqrt n\sum_{j=0}^p a_j=0$，每个余项 $\sqrt{n+j}-\sqrt n\to0$。第三题用 $\sqrt{1+t}-1=t/(\sqrt{1+t}+1)$，其和趋于 $\tfrac12\cdot\tfrac12=\tfrac14$。指数比值趋于 $1$；对数比值趋于 $3$。
\item \textbf{平均值型结论：}Cesàro 定理给出前 $n$ 项算术平均的极限；对正项取对数可得几何平均的极限。Stolz 公式由差分比的加权平均得到；奇数平方和 $\sum_{k=1}^n(2k-1)^2=n(4n^2-1)/3$，故所求极限为 $4/3$。
\item \textbf{$1、2、5、6$ 班定义题：}“存在无限多个 $\varepsilon>0$”及“对每个 $\varepsilon>1$”不足以刻画收敛，例如常数列 $a_n=a+1/2$。对每个 $0<\varepsilon<1$、对每个 $k$ 的 $1/k$，以及任意两侧邻域条件，均与 $a_n\to a$ 等价。整数值收敛数列最终恒定；有界数列的夹逼须两端趋于同一极限，单有两端间距趋零不保证中间列收敛。
\item \textbf{选做：}三列两两之差每次乘 $-1/2$，而总和始终为 $a+b+c$，故三列均趋于 $(a+b+c)/3$。Fibonacci 相邻项比值的极限是正根 $r$ 满足 $r=1/(1+r)$，即 $(\sqrt5-1)/2$。Toeplitz 定理将有限项与尾项分开估计：尾项由收敛性控制，有限项系数逐列趋零。
\end{itemize}

% ANSWERS: f4
\begin{itemize}
\item \textbf{教材章节题号：}参见教材。
\item \textbf{第 3 次作业中重复的极限与 Cesàro 题：}参见前一次的对应补充答案。
\item \textbf{共同子列：}由 $[a,b]$ 上的有界性，先取 $x_{n_k}\to\alpha$；因 $x_n-y_n\to0$，同一下标的 $y_{n_k}\to\alpha$。
\item \textbf{单调数列与子列：}单调列若有有限收敛子列，则不能向 $\pm\infty$ 发散；由单调收敛定理，整列收敛到同一极限。
\item \textbf{阶乘极限：}令 $a_n=2^nn!/n^n$，则 $a_{n+1}/a_n=2(n/(n+1))^n\to2/e<1$，故 $a_n\to0$。
\item \textbf{柯西列：}给定 $\varepsilon$，选足够靠后的收敛子列项并用柯西性质估计所有后续项，原列即收敛。$\sin n$ 不满足柯西准则；若它收敛到 $L$，由 $\sin(n+1)-\sin(n-1)=2\sin1\cos n$ 得 $\cos n\to0$，而递推式又给 $L=L\cos1$，故 $L=0$，与 $\sin^2n+\cos^2n=1$ 矛盾。
\item \textbf{条件核对：}原稿“数列有收敛子列当且仅当该数列不是无穷大数列”须明确“无穷大数列”指 $|a_n|\to\infty$；按此定义，由 Bolzano--Weierstrass 定理结论成立。
\end{itemize}

% ANSWERS: f5
\begin{itemize}
\item \textbf{教材章节题号：}参见教材。
\item \textbf{前次重复题：}共同收敛子列、单调列、阶乘极限及柯西判别等见第 4 次作业的补充答案。
\item \textbf{函数极限：}单调函数的一侧极限等于相应区间上的上确界或下确界；若所有趋于 $+\infty$ 的严格增列上均有 $f(x_n)\to A$，反设函数极限不为 $A$，可递归选出违背条件的增列。
\item \textbf{周期函数差：}连续性并非题设条件；若两个周期函数的差趋零，可用周期平移及同余逼近论证两函数相同，证明中应区分两周期可公度与不可公度的情形。
\item \textbf{选做题：}由 $0<e-\sum_{k=0}^n1/k!<1/(n\,n!)$ 可证 $e$ 无理。二进制序列对应的级数可用几何级数估计证明收敛；端点的两种表示需要单独合并处理。
\end{itemize}

% ANSWERS: f6
\begin{itemize}
\item \textbf{教材章节题号：}参见教材。
\item \textbf{局部极限与有界性：}每点有有限极限意味着该点附近函数有界；有限个邻域覆盖紧区间 $[a,b]$，故 $f$ 全局有界。
\item \textbf{大 $O$ 与小 $o$：}若 $|g(x)|\le Cx^2$，则 $|g(x)/x|\le C|x|\to0$。$|\cos x-1|\le x^2/2$ 给出 $\cos x=1+O(x^2)$。
\item \textbf{极限：}因 $\tan5x\sim5x$、$\sin3x\sim3x$ 且 $1-\cos x=O(x^2)$，所求极限为 $3/5$。
\item \textbf{Cauchy 列保持性质：}若 $f$ 在某点不连续，可取趋于该点的数列并与常数子列交错，构造 Cauchy 输入而非 Cauchy 输出，与题设矛盾。
\end{itemize}

% ANSWERS: f7
\begin{itemize}
\item \textbf{多项式及连续函数：}首一偶次多项式两端都趋 $+\infty$，可以没有零点；若常数项为负，由 $f(0)<0$ 和两端正值，用介值定理得至少两个实零点。若 $f$ 在 $(a,b)$ 两端趋 $+\infty$，把两端外侧截去后在紧区间上取最小值。
\item \textbf{平移中值：}令 $g(x)=f(x)-f(x+a)$，$x\in[0,a]$。由于 $g(0)=-g(a)$，介值定理给出零点。
\item \textbf{唯一实根：}$h(x)=x^3+2x-1$ 严格递增，且 $h(0)<0<h(1)$，故恰有一个零点。
\item \textbf{周期与一致连续：}在一周期的紧区间上用 Heine--Cantor 定理取共同 $\delta$；通过周期平移把任意两点移回有限个相邻周期即可。
\item \textbf{选做题：}开集与闭集的等价刻画可利用补集和序列收敛；连续模 $\omega(\delta)\to0$ 当且仅当函数一致连续（定义域按原题）。
\end{itemize}

% ANSWERS: f8
\begin{itemize}
\item \textbf{一致连续：}和、差保持一致连续；乘积须额外有界性。$\sin x^2$ 在全实轴不一致连续，取 $x_n=\sqrt{2\pi n+\pi/2}$ 与 $y_n=\sqrt{2\pi n+3\pi/2}$，则 $|x_n-y_n|\to0$ 而函数值相差 $2$。
\item \textbf{可导性：}若 $f$ 在 $a$ 连续且 $f(a)\ne0$，并且 $f^2$ 在 $a$ 可导，则 $f'(a)=(f^2)'(a)/(2f(a))$。若 $|f|$ 在零点可导，则其导数只能为 $0$，从 $|f(x)|=o(|x-a|)$ 得 $f'(a)=0$；非零点局部符号不变。
\item \textbf{例子与运算：}若偶函数在 $0$ 可导，则 $f'(0)=0$。$f(x)=(e^x-1)\prod_{k=2}^n(e^{kx}-k)$，故 $f'(0)=\prod_{k=2}^n(1-k)=(-1)^{n-1}(n-1)!$。
\end{itemize}

% ANSWERS: f9
\begin{itemize}
\item \textbf{教材第 5 章各节及总练题号：}参见教材。
\item \textbf{偶函数：}若 $f$ 在 $0$ 可导，由 $f(h)=f(-h)$，两侧差商互为相反数，故 $f'(0)=0$。
\item \textbf{高阶导数：}$2x/(1-x^2)=1/(1-x)-1/(1+x)$，因此 $f^{(n)}(x)=n!\bigl[(1-x)^{-n-1}-(-1)^n(1+x)^{-n-1}\bigr]$（$|x|<1$）。
\item \textbf{多变量复合微分：}按原题各小问先写全微分，再对 $u(x),v(x),w(x)$ 应用链式法则。
\end{itemize}

% ANSWERS: f10
\begin{itemize}
\item \textbf{教材 6.1、6.2 节题号：}参见教材。
\item \textbf{非线性函数：}若任意 $x_1<x_2$ 的割线斜率都相同，则函数线性；故非线性时存在两段割线斜率不同，再由中值定理得到两点导数不同。
\item \textbf{积分/中值题：}各题可先构造题设等式两边之差，再用 Rolle 定理、Lagrange 中值定理或 Cauchy 中值定理；涉及 $f(a)=f(b)=0$ 的二阶式，可在 $a,x,b$ 三点插值并对插值余项使用 Rolle 定理。
\item \textbf{条件校核：}涉及开区间端点极限时，须先把函数连续延拓到闭区间，才能直接应用中值定理。
\end{itemize}

% ANSWERS: f11
\begin{itemize}
\item \textbf{教材 6.2、6.3 节及总练题号：}参见教材。
\item \textbf{反正弦展开：}$\arcsin x=\sum_{k=0}^{\infty}\dfrac{\binom{2k}{k}}{4^k(2k+1)}x^{2k+1}$，$|x|<1$；可由 $(\arcsin x)'=(1-x^2)^{-1/2}$ 的二项式级数逐项积分得到。
\item \textbf{根式极限：}二阶 Taylor 展开给出 $\sqrt{x+1}+\sqrt{x-1}-2\sqrt x\sim-1/(4x^{3/2})$，故极限为 $-1/4$。
\item \textbf{高阶导数：}$x^2 2^x=x^2e^{x\ln2}$ 的 $x^n$ 系数为 $(\ln2)^{n-2}/(n-2)!$，故 $f^{(n)}(0)=n(n-1)(\ln2)^{n-2}$（$n\ge2$）；$n=0,1$ 时均为 $0$。
\end{itemize}

% ANSWERS: s1
\begin{itemize}
\item \textbf{教材 8.3 节题号：}参见教材。
\item \textbf{代换法：}$\int e^{\sqrt[3]x}\,dx$ 令 $t=\sqrt[3]x$，得 $3e^t(t^2-2t+2)+C$。其他含根式或三角函数的不定积分，先作对应代换再分部积分。
\item \textbf{递推积分：}$\int\cos^n x\,dx=\sin x\cos^{n-1}x/n+(n-1)\int\cos^{n-2}x\,dx/n$；$\int\sec^n x\,dx=\tan x\sec^{n-2}x/(n-1)+(n-2)\int\sec^{n-2}x\,dx/(n-1)$（$n\ge2$）。
\item \textbf{分式积分：}把分子写成分母导数与常数项的线性组合；对 $1/(1+x^3)$ 先分解 $1+x^3=(x+1)(x^2-x+1)$，再作部分分式。
\end{itemize}

% ANSWERS: s2
\begin{itemize}
\item \textbf{教材 9.3、9.6 节题号：}参见教材。
\item \textbf{Riemann 可积函数的线性与乘积：}用上下和判别或振幅和判别，先证明有限线性组合可积；乘积可由 $fg=((f+g)^2-f^2-g^2)/2$ 化为平方的可积性。
\item \textbf{局部上确界条件：}若每个子区间都有 $\sup_\Delta f\ge\sigma$，则上和至少为 $\sigma(b-a)$；可积时下、上和有共同极限，因此 $\int_a^b f\ge\sigma(b-a)$。
\item \textbf{每点极限为零：}每点附近可使函数值足够小（点值本身可能不同）；利用有限覆盖只能控制去掉有限个点后的积分，结论应以原题最终所求量为准。
\end{itemize}

% ANSWERS: s3
\begin{itemize}
\item \textbf{教材 9.4 节题号：}参见教材。
\item \textbf{积分中值定理：}连续 $f$ 在 $[a,b]$ 上取极值 $m,M$，故 $m(b-a)\le\int_a^b f\le M(b-a)$，由介值定理存在 $\xi$ 使积分等于 $f(\xi)(b-a)$；有非负权 $g$ 时同理夹住 $\int fg$。
\item \textbf{积分型 Cauchy--Schwarz：}对任意实数 $t$ 有 $\int_a^b(tf-g)^2\ge0$；判别式非正，得到 $(\int fg)^2\le(\int f^2)(\int g^2)$。
\item \textbf{单调函数积分：}可用上下和及微积分基本定理处理原题中的单调性结论；若要求严格不等式，需核对题设是否排除了常数函数。
\end{itemize}

% ANSWERS: s4
\begin{itemize}
\item \textbf{零积分：}若存在 $x_0$ 使 $f(x_0)>0$，连续性给出一段正长度区间，其积分严格正，与非负函数总积分为零矛盾。
\item \textbf{积分交换：}$\int_0^\pi\int_0^x\frac{\sin t}{\pi-t}\,dt\,dx=\int_0^\pi\sin t\,dt=2$；用积分区域交换次序。
\item \textbf{参数积分：}对原题 $B(m,n)=\int_0^1x^{m-1}(1-x)^{n-1}\,dx$，反复分部积分得 $(m-1)!(n-1)!/(m+n-1)!$。
\item \textbf{其余证明：}优先由微积分基本定理构造原函数；变上限积分的单调性由差值 $\int_{x_1}^{x_2}f$ 的符号判断。
\end{itemize}

% ANSWERS: s5
\begin{itemize}
\item \textbf{取整函数：}每个 $[k,k+1]$ 上 $x-[x]=x-k$，积分为 $1/2$，故 $\int_0^n(x-[x])\,dx=n/2$。
\item \textbf{对称积分：}$\int_{-1}^1\!\frac{dx}{(e^x+1)(1+x^2)}=\pi/4$，因 $1/(e^x+1)+1/(e^{-x}+1)=1$。$\int_0^{\pi/4}\!\frac{\sin x}{\sin x+\cos x}\,dx=\pi/8-\tfrac14\ln2$。
\item \textbf{对数积分：}$\int_0^1(\ln x)^n\,dx=(-1)^n n!$，用 $x=e^{-t}$ 化为 Gamma 积分。
\item \textbf{收敛积分与被积函数极限：}若 $\int_a^\infty f$ 收敛且 $f(x)$ 的有限极限存在，则该极限必为零；若 $f$ 单调，存在扩展实数极限，再结合收敛性可得 $f(x)\to0$。
\end{itemize}

% ANSWERS: s6
\begin{itemize}
\item \textbf{对数幂积分：}原题下限 $2$ 使 $\ln\ln x$ 在积分区间内经过 $0$，且对非整数 $p$ 不一定有实数定义；本册把下限改为 $e^e$。此时令 $u=\ln\ln x$，所得积分当且仅当 $p>1$ 收敛。
\item \textbf{反余弦积分：}$u=\arccos x$ 给出 $\int_{-1}^1\frac{\arccos x}{\sqrt{1-x^2}}\,dx=\int_0^\pi u\,du=\pi^2/2$。
\item \textbf{正项积分分组：}非负 $f$ 在相邻区间 $[A_n,A_{n+1}]$ 上的积分为非负数；反常积分收敛当且仅当这些积分组成的正项级数收敛。
\item \textbf{一致连续与衰减：}若 $f$ 一致连续而 $\int_a^\infty f$ 收敛，反设 $|f(x_n)|\ge\varepsilon$，一致连续性给出一系列固定宽度的区间，使积分的 Cauchy 条件矛盾；须用适当的单侧同号区间。
\end{itemize}

% ANSWERS: s7
\begin{itemize}
\item \textbf{比较判别法：}$\int_0^\infty(1+x^5)^{-1/2}\,dx$ 收敛；$\int_1^\infty x^pe^{-x}\,dx$ 对 $p\ge0$ 收敛；$\int_1^\infty(\arctan x)x^{-p}\,dx$ 当且仅当 $p>1$ 收敛；$\int_1^\infty r^{-\ln x}\,dx=\int_1^\infty x^{-\ln r}\,dx$ 当且仅当 $r>e$ 收敛。
\item \textbf{端点局部展开：}$\int_0^1x^{-p}\sin(1/x)\,dx$ 令 $t=1/x$，化为 $\int_1^\infty t^{p-2}\sin t\,dt$；$p<1$ 时绝对收敛，$1\le p<2$ 时条件收敛。$\int_0^1x^{-p}\sin x\,dx$ 当且仅当 $p<2$ 收敛。
\item \textbf{其他典型比较：}$1-\cos(2/x)\sim2/x^2$，故相应无穷积分收敛；$\cos(1/x)-1\sim-1/(2x^2)$，亦收敛。$|\ln x|^p$ 在 $0$ 处当且仅当 $p>-1$ 可积。
\item \textbf{其余比较题：}第 7 小题因被积函数在无穷远处与 $\sin^2x/x^{2p}$ 同阶，在题设 $2p>1$ 下收敛。第 9 小题被积函数与常数倍 $x^{-1/2}$ 同阶，发散。第 10 小题与常数倍 $x^{-1-p/2}$ 同阶，当且仅当 $p>0$ 收敛。第 13--16 小题分别在端点与 $(1-x)^{-1/2}$、$x^{-1/2}$、$x^{-2/3}$、$x^{-1/2}$ 同阶，均收敛。
\item \textbf{剩余参数题：}第 18 小题 $x^{-p}\ln\cos(1/x)\sim-\tfrac12x^{-p-2}$，当且仅当 $p>-1$ 收敛。第 19 小题还须保证 $x^2+p$ 在积分区间无零点；在 $p>-1$ 下，被积函数主项为 $(1-p)/x$，故当且仅当 $p=1$ 收敛。
\item \textbf{振荡积分：}$\int_0^\infty x^p\sin x\,dx$ 对题设 $p>0$ 发散；分部积分后的边界项不能趋零。$\int_0^\infty(x-\sin x)x^{-3}\,dx$ 两端比较均可积。
\end{itemize}

% ANSWERS: s8
\begin{itemize}
\item \textbf{第 1 题：}按原稿给出的分解，$\int_1^\infty\sin x/x^p\,dx$ 对 $p>0$ 由 Dirichlet 判别收敛；余下的非负积分与 $\int_1^\infty\sin^2x/x^{2p}\,dx$ 同敛散，阈值为 $p=1/2$。绝对收敛阈值为 $p>1$。零点附近被积函数局部可积。
\item \textbf{Gamma/Beta 型积分：}$\int_0^\infty x^{s-1}e^{-x}\,dx$ 当且仅当 $s>0$ 收敛；$\int_0^1x^{p-1}(1-x)^{q-1}\,dx$ 当且仅当 $p,q>0$ 收敛。
\item \textbf{反例：}取 $f(x)=\sin(x^2)$，则 Fresnel 型积分 $\int_1^\infty f$ 收敛，但 $\int_1^\infty f^2$ 发散。要展示条件收敛而非绝对收敛，可取 $f(x)=\sin x/x$。
\item \textbf{变量代换与条件修正：}原稿的 $F(1)=0$，且积分区间内 $F$ 会穿过 $1$，原式的对数分母有奇点。本册改从满足 $F(c)=e$ 的 $c$ 起积分。令 $u=F(x)$，$I=\int_e^\infty du/(u\ln u)$ 发散，$J=\int_e^\infty du/(u\ln^2u)$ 收敛。
\item \textbf{凝聚判别：}正项递减 $a_n$ 满足 $2^ka_{2^{k+1}}\le\sum_{n=2^k}^{2^{k+1}-1}a_n\le2^ka_{2^k}$，于是 $\sum a_n$ 与 $\sum2^ka_{2^k}$ 同敛散。由此 $\sum n^{-p}$ 当且仅当 $p>1$ 收敛，$\sum_{n\ge2}(n\ln^p n)^{-1}$ 当且仅当 $p>1$ 收敛。
\end{itemize}

% ANSWERS: s9
\begin{itemize}
\item \textbf{正项级数：}比较、积分及 Cauchy 凝聚判别应先核对正项和单调条件；若题目包含参数，需分别讨论边界参数值。
\item \textbf{广义二项式系数：}对非整数 $a>0$，比值 $|\binom a{n+1}/\binom a n|=|a-n|/(n+1)=1-(a+1)/(n+1)$ 当 $n$ 足够大；由 Raabe 判别法得 $\sum_n|\binom a n|$ 收敛。正整数 $a$ 时仅有限项非零。
\item \textbf{收敛项与级数：}若 $a_n\to a\ne0$，需先检查所论级数的通项是否趋零；通项不趋零的级数必发散。
\item \textbf{分块符号级数：}按 $k^2\le n<(k+1)^2$ 分组，块的绝对值约为 $2/k$，符号 $(-1)^k$，故 $\sum(-1)^{[\sqrt n]}/n$ 由交错判别收敛，但其绝对值级数为调和级数，发散。
\end{itemize}

% ANSWERS: s10
\begin{itemize}
\item \textbf{阶乘控制：}在 $1/2\le|x|\le2$ 上，$|x^n+x^{-n}|\le2^{n+1}$；$\sum 2^{n+1}n^2/\sqrt{n!}$ 由比值判别收敛，故原函数项级数一致收敛。
\item \textbf{单调函数与端点：}若 $u_n$ 在 $[a,b]$ 单调，则 $|u_n(x)|\le\max(|u_n(a)|,|u_n(b)|)$；端点绝对收敛给出一致收敛的 Weierstrass 判别。
\item \textbf{一致趋零选子列：}递归选 $n_k$ 使 $\sup_{x\in I}|f_{n_k}(x)|<2^{-k}$，再用 M 判别法。
\item \textbf{Dirichlet 级数：}$\sum_{n\ge1}\sin(nx)/n$ 在任意闭子区间 $[\delta,2\pi-\delta]$ 上一致收敛，在整个 $(0,2\pi)$ 上不一致收敛；前者用三角部分和一致有界，后者利用靠近端点时部分和的非一致 Cauchy 性。
\item \textbf{积分与求和：}$S(t)=\sum_{n\ge1}ne^{-nt}$ 在 $[\ln2,\ln3]$ 一致收敛，因此 $\int_{\ln2}^{\ln3}S(t)\,dt=\sum_{n\ge1}(2^{-n}-3^{-n})=1-1/2=1/2$。
\end{itemize}
